\begin{tikzpicture}[x=1cm,y=1cm,>=Latex,font=\small,text=navy]
\path[use as bounding box] (-.1,.6) rectangle (16.4,-19.7);
\node[anchor=center,font=\sffamily\bfseries\small,text=violet] at (8.1,.3) {MAPA EPISTÉMICO · CÓMO SE SOSTIENE CADA RESULTADO};
\node[draw=navy!35,fill=navy!3,rounded corners=3pt,text width=15.8cm,align=center,inner sep=8pt] at (8.1,-.7) {\textbf{Premisa física y operador de observación}\\[3pt]Plantilla efectiva explícita; calibración de masa; excentricidad; predicciones propias de cada teoría.};
\node[anchor=west,font=\sffamily\bfseries\footnotesize,text=violet] at (0.0,-1.7) {RESULTADO};
\node[anchor=west,font=\sffamily\bfseries\footnotesize,text=violet] at (5.55,-1.7) {CÁLCULO Y EVIDENCIA};
\node[anchor=west,font=\sffamily\bfseries\footnotesize,text=violet] at (11.1,-1.7) {APLICACIÓN Y CONDICIONES};
\node[draw=navy!55,fill=navy!3,rounded corners=3pt,text width=4.65cm,minimum height=1.95cm,align=left,inner sep=7pt] (physics0) at (2.45,-3.0) {\textbf{Física y calibración}\\[6pt]$a_r=-\mu[1+\alpha(1+x)e^{-x}]/r^2$\\[5pt]Gradiente, calibración de masa y excentricidad finita.};
\node[draw=navy!55,fill=navy!3,rounded corners=3pt,text width=4.65cm,minimum height=1.95cm,align=left,inner sep=7pt] (physics1) at (8.0,-3.0) {\textbf{Tres rutas de cálculo}\\[6pt]Bessel, cuadratura e integración\\[5pt]Error semilatus: $-67.63\%$ ($e=0.82697$, $x=1$).};
\node[draw=navy!55,fill=navy!3,rounded corners=3pt,text width=4.65cm,minimum height=1.95cm,align=left,inner sep=7pt] (physics2) at (13.55,-3.0) {\textbf{Calibrar la respuesta}\\[6pt]Plantilla fenomenológica; propagación luminosa con hipótesis métricas explícitas.};
\draw[->,navy!60,thick] (physics0.east)--(physics1.west);
\draw[->,navy!60,thick] (physics1.east)--(physics2.west);
\node[draw=violet!55,fill=violet!3,rounded corners=3pt,text width=4.65cm,minimum height=1.95cm,align=left,inner sep=7pt] (gp0) at (2.45,-5.65) {\textbf{Ley conjunta GP}\\[6pt]$Db=Db^{\prime}$; $DKD^T=DK^{\prime}D^T$\\[5pt]La ley posterior se determina por sus contrastes.};
\node[draw=violet!55,fill=violet!3,rounded corners=3pt,text width=4.65cm,minimum height=1.95cm,align=left,inner sep=7pt] (gp1) at (8.0,-5.65) {\textbf{Distancia entre leyes}\\[6pt]$D_{\rm KL}=1.402585$ nats\\[5pt]Control sintético de covarianza; soportes singulares explícitos.};
\node[draw=violet!55,fill=violet!3,rounded corners=3pt,text width=4.65cm,minimum height=1.95cm,align=left,inner sep=7pt] (gp2) at (13.55,-5.65) {\textbf{Auditar emuladores}\\[6pt]Criterio en malla finita; estabilidad posterior y conectividad FractalGo.};
\draw[->,violet!60,thick] (gp0.east)--(gp1.west);
\draw[->,violet!60,thick] (gp1.east)--(gp2.west);
\node[draw=plum!55,fill=plum!3,rounded corners=3pt,text width=4.65cm,minimum height=1.95cm,align=left,inner sep=7pt] (cassini0) at (2.45,-8.3) {\textbf{Jerarquía Cassini}\\[6pt]$y=A s+X\beta+\epsilon$\\[5pt]1696 épocas reconstruidas; cuatro covarianzas temporales.};
\node[draw=plum!55,fill=plum!3,rounded corners=3pt,text width=4.65cm,minimum height=1.95cm,align=left,inner sep=7pt] (cassini1) at (8.0,-8.3) {\textbf{Localizar la discrepancia}\\[6pt]Pasada 12: $0/1024$ excedencias\\[5pt]Discrepancia temporal localizada por pasada.};
\node[draw=plum!55,fill=plum!3,rounded corners=3pt,text width=4.65cm,minimum height=1.95cm,align=left,inner sep=7pt] (cassini2) at (13.55,-8.3) {\textbf{Evaluar el modelo temporal}\\[6pt]Mezcla exploratoria sobre datos reconstruidos; amplitud óptica distinta de $(\alpha,\lambda)$.};
\draw[->,plum!60,thick] (cassini0.east)--(cassini1.west);
\draw[->,plum!60,thick] (cassini1.east)--(cassini2.west);
\node[draw=cyan!55,fill=cyan!3,rounded corners=3pt,text width=4.65cm,minimum height=1.95cm,align=left,inner sep=7pt] (information0) at (2.45,-10.95) {\textbf{Información y decisión}\\[6pt]$K_{\pi,L}=1-R_{\pi,L}(E)/R_{\pi,L}(\varnothing)$\\[5pt]Retención relativa a un objeto, prior y pérdida declarados.};
\node[draw=cyan!55,fill=cyan!3,rounded corners=3pt,text width=4.65cm,minimum height=1.95cm,align=left,inner sep=7pt] (information1) at (8.0,-10.95) {\textbf{Fisher, Rényi y riesgo}\\[6pt]$J_{\rm mix}=0.345509<0.555556$\\[5pt]Fisher y riesgo resuelven el control de mezcla.};
\node[draw=cyan!55,fill=cyan!3,rounded corners=3pt,text width=4.65cm,minimum height=1.95cm,align=left,inner sep=7pt] (information2) at (13.55,-10.95) {\textbf{Comparar transformaciones}\\[6pt]Objetivo, prior y pérdida declarados; riesgo basal positivo y finito.};
\draw[->,cyan!60,thick] (information0.east)--(information1.west);
\draw[->,cyan!60,thick] (information1.east)--(information2.west);
\node[draw=navy!55,fill=navy!3,rounded corners=3pt,text width=4.65cm,minimum height=1.95cm,align=left,inner sep=7pt] (cases0) at (2.45,-13.6) {\textbf{Gaia y QUMOND}\\[6pt]$C_{\rm shared};\ Q_2$\\[5pt]Covarianza Gaia por tránsito; respuesta QUMOND al campo externo.};
\node[draw=navy!55,fill=navy!3,rounded corners=3pt,text width=4.65cm,minimum height=1.95cm,align=left,inner sep=7pt] (cases1) at (8.0,-13.6) {\textbf{Convergencia y procedencia}\\[6pt]$Q_2=3.38677\times10^{-26}$ $\mathrm{s}^{-2}$\\[5pt]Registros ESA y cuadratura del modelo de campo externo.};
\node[draw=navy!55,fill=navy!3,rounded corners=3pt,text width=4.65cm,minimum height=1.95cm,align=left,inner sep=7pt] (cases2) at (13.55,-13.6) {\textbf{Relacionar geometría y física}\\[6pt]Gaia a diseño fijo; cuadrupolo teórico con campo e interpolación declarados.};
\draw[->,navy!60,thick] (cases0.east)--(cases1.west);
\draw[->,navy!60,thick] (cases1.east)--(cases2.west);
\node[draw=violet!55,fill=violet!3,rounded corners=3pt,text width=4.65cm,minimum height=1.95cm,align=left,inner sep=7pt] (dcae0) at (2.45,-16.25) {\textbf{DCAE y COSMOS}\\[6pt]$I_j^\uparrow=[1-\mathcal R_j]_+$\\[5pt]v0.6: componente, protocolo y $\Delta_j$ firmado.};
\node[draw=violet!55,fill=violet!3,rounded corners=3pt,text width=4.65cm,minimum height=1.95cm,align=left,inner sep=7pt] (dcae1) at (8.0,-16.25) {\textbf{COSMOS: control adverso}\\[6pt]$\Delta_{\rm soporte}=-0.0092143$\\[5pt]COSMOS: el orden cambia con tolerancia y observable.};
\node[draw=violet!55,fill=violet!3,rounded corners=3pt,text width=4.65cm,minimum height=1.95cm,align=left,inner sep=7pt] (dcae2) at (13.55,-16.25) {\textbf{Describir el acceso}\\[6pt]Descriptor de soporte en tablas; medición física DCAE pendiente.};
\draw[->,violet!60,thick] (dcae0.east)--(dcae1.west);
\draw[->,violet!60,thick] (dcae1.east)--(dcae2.west);
\node[draw=violet!40,fill=violet!5,rounded corners=2pt,text width=15.8cm,align=center,inner sep=7pt] at (8.1,-18.5) {\textbf{Identificabilidad + ley conjunta + adecuación predictiva}\\[4pt]GRAVITAS + FractalGo: reproducir los cálculos, inspeccionar la evidencia y extender los contrastes.};
\node[anchor=west,font=\tiny,text=navy!65] at (0,-19.5) {Mismo modelo de contenidos que el grafo interactivo. SHA-256: ab0cb2ad30b4e0d8.};
\end{tikzpicture}
